\documentclass[12pt]{article}
\usepackage[utf8]{inputenc}
\usepackage{graphicx}
\usepackage{amsmath}
\usepackage{amssymb}
\usepackage{mathrsfs}
\usepackage{pgfornament}
\usepackage[many]{tcolorbox}
\usepackage[margin = 0.5in]{geometry}
\usepackage{collectbox}
\usepackage{mdframed}
\usepackage{xcolor}

\newcommand\textbox[1]{%
  \parbox{.333\textwidth}{#1}%
}

\linespread{1.5}

\makeatletter
\newcommand{\mybox}{%
    \collectbox{%
        \setlength{\fboxsep}{2pt}%
        \fbox{\BOXCONTENT}%
    }%
}
\makeatother

\usepackage{listings}
\usepackage{color}

\definecolor{dkgreen}{rgb}{0,0.6,0}
\definecolor{gray}{rgb}{0.5,0.5,0.5}
\definecolor{mauve}{rgb}{0.58,0,0.82}

\lstset{frame=tb,
  language=R,
  aboveskip=3mm,
  belowskip=3mm,
  showstringspaces=false,
  columns=flexible,
  basicstyle={\small\ttfamily},
  numbers=none,
  numberstyle=\tiny\color{gray},
  keywordstyle=\color{blue},
  commentstyle=\color{dkgreen},
  stringstyle=\color{mauve},
  breaklines=true,
  breakatwhitespace=true,
  tabsize=3
}
\title{STS 532 HW 4}
\begin{document}

%%
%% Title 
%%
\begin{center}
\vspace{-0.6cm}
\hrulefill \\
\vspace{-0.7cm}
\hrulefill

\vspace{-0.5cm}
\begin{flushleft}
\noindent\textbox{\pgfornament[width = 9mm, color = black]{39}\hfill}\textbox{\hfil \textbf{\begin{large}Mathematical Statistics\end{large}}\hfil}\textbox{\hfill \pgfornament[width = 9mm, color = black]{40}}
\end{flushleft}
\vspace{-5mm}
\begin{large}
\textbf{Homework 4} \\
\end{large}

\vspace{5mm}

\begin{flushleft}
\noindent\textbox{\pgfornament[width = 9mm, color = black, symmetry = h]{39}\hfill}\textbox{\hfil \textbf{\begin{large}Nolan R. H. Gagnon\end{large}}\hfill}\textbox{\hfill \pgfornament[width = 9mm, color = black, symmetry = h]{40}}
\end{flushleft}

\vspace{-0.65cm}
\hrulefill \\
\vspace{-0.7cm}
\hrulefill
\end{center}

\begin{center}
\pgfornament[width = 75mm,
             color = black]{89}
\end{center}

\vspace{5mm}

\noindent\mybox{\colorbox{lightgray}{\textbf{Wackerly 10.126}}}\hspace{3mm}\hrulefill 

\noindent Suppose that $X_1, X_2, \ldots, X_{n_1}$, $Y_1, Y_2, \ldots, Y_{n_2}$, and $W_1, W_2, \ldots, W_{n_3}$ are independent random samples from normal distributions with respective unknown means $\mu_1, \mu_2,$ and $\mu_3$ and common variances $\sigma_1^2 = \sigma_2^2=\sigma_3^2=\sigma_2$.  Suppose that we want to estimate a linear function of the means: $\theta = a_1\mu_1 + a_2\mu_2 + a_3\mu_3$. Because the maximum-likelihood estimator (MLE) of a function of parameters is the function of the MLEs of the parameters, the MLE of $\theta$ is $\hat{\theta} = a_1\bar{X}+a_2\bar{Y}+a_3\bar{W}.$
\begin{flushleft}
\textbf{(a)} What is the standard error of the estimator $\hat{\theta}$? \\
\textbf{(b)} What is the distribution of $\hat{\theta}$? \\

\textbf{(c)} If the sample variances are given by $S_1^2, S_2^2,$ and $S_3^2$, respectively, consider $$S_p^2 = \frac{(n_1-1)S_1^2+(n_2-1)S_2^2+(n_3-1)S_3^2}{n_1+n_2+n_3 - 3}.$$
\hspace{5mm} \textbf{(i)} What is the distribution of $(n_1+n_2+n_3-3)S_p^2/\sigma^2?$ \\
\hspace{5mm} \textbf{(ii)} What is the distribution of $$T=\frac{\hat{\theta}-\theta}{S_p\sqrt{\frac{a_1^2}{n_1}+\frac{a_2^2}{n_2}+\frac{a_3^2}{n_3}}}?$$ \\
\textbf{(d)} Give a confidence interval for $\theta$ with confidence coefficient $1-\alpha$. \\
\textbf{(e)} Develop a test for $H_0: \theta = \theta_0$ versus $H_1: \theta\neq\theta_0$.
\end{flushleft}

\noindent\mybox{\colorbox{lightgray}{\textbf{Solution}}} \hspace{3mm}\hrulefill

\noindent \textcolor{blue}{\textbf{(a)}} Since the samples are independent of each other, the variance of $\hat{\theta}$ is
\begin{align}
\text{Var}\left(\hat{\theta}\right) &= \text{Var}\left(a_1\bar{X}+a_2\bar{Y}+a_3\bar{W}\right) \\
&= a_1^2\text{Var}\left(\bar{X}\right)+a_2^2\text{Var}\left(\bar{Y}\right)+a_3^2\text{Var}\left(\bar{W}\right) \\
&= a_1^2\frac{\sigma^2}{n_1} + a_2^2\frac{\sigma^2}{n_2} + a_3\frac{\sigma^2}{n_3} \\
&= \sigma^2\left(\frac{a_1^2}{n_1} + \frac{a_2^2}{n_2} + \frac{a_3^2}{n_3}\right).
\end{align}
The standard error of $\hat{\theta}$ is
\begin{align}
\text{SE}\left(\hat{\theta}\right) &= \sqrt{\text{Var}\left(\hat{\theta}\right)} \\
&= \sigma\sqrt{\frac{a_1^2}{n_1} + \frac{a_2^2}{n_2} + \frac{a_3^2}{n_3}}.
\end{align}

\vspace{3mm}

\noindent \textcolor{blue}{\textbf{(b)}} By Theorem 5.3.1 (Casella and Berger $\S$ 5.3),
\begin{align}
\bar{X} &\sim n\left(\mu_1, \frac{\sigma^2}{n_1} \right) \\
\bar{Y} &\sim n\left(\mu_2, \frac{\sigma^2}{n_2} \right) \\
\bar{W} &\sim n\left(\mu_3, \frac{\sigma^2}{n_3} \right).
\end{align}
Using the method of moment-generating functions, it is easy to see that
\begin{align}
a_1\bar{X} &\sim n\left(a_1\mu_1, \frac{a_1^2\sigma^2}{n_1}\right) \\
a_2\bar{Y} &\sim n\left(a_2\mu_2, \frac{a_2^2\sigma^2}{n_1}\right) \\
a_3\bar{W} &\sim n\left(a_3\mu_3, \frac{a_3^2\sigma^2}{n_1}\right).
\end{align}
Thus $\hat{\theta}$ is normal with mean $a_1\mu_1 + a_2\mu_2 + a_3\mu_3 = \theta$, and variance $\sigma^2\left(\frac{a_1^2}{n_1} + \frac{a_2^2}{n_2} + \frac{a_3^2}{n_3}\right)$.

\vspace{3mm}

\noindent \textcolor{blue}{\textbf{(c)}} \textcolor{magenta}{\textbf{(i)}} Notice that we can write
\begin{align}
(n_1+n_2+n_3-3)S_p^2/\sigma^2 &= \frac{(n_1-1)S_1^2}{\sigma^2} + \frac{(n_2-1)S_2^2}{\sigma^2} + \frac{(n_3-1)S_3^2}{\sigma^2}.
\end{align}
By Theorem 5.3.1 (Casella and Berger $\S$ 5.3),
\begin{align}
\frac{(n_1-1)S_1^2}{\sigma^2} &\sim \chi^2(n_1-1) \\
\frac{(n_2-1)S_2^2}{\sigma^2} &\sim \chi^2(n_2-1) \\
\frac{(n_3-1)S_3^2}{\sigma^2} &\sim \chi^2(n_3-1)
\end{align}
Note that the above variables are independent. Now, a sum of $k$ independent chi-squared random variables (with degrees of freedom $p_i$) is again a chi-squared random variable with $\sum\limits_{i=1}^k p_i$ degrees of freedom. Thus, $$(n_1+n_2+n_3-3)S_p^2/\sigma^2 \sim \chi^2\left(n_1+n_2+n_3-3\right).$$

\vspace{3mm}

\noindent\textcolor{magenta}{\textbf{(ii)}} We can write
\begin{align}
T &= \frac{\hat{\theta}-\theta}{S_p\sqrt{\frac{a_1^2}{n_1}+\frac{a_2^2}{n_2}+\frac{a_3^2}{n_3}}} \\
&= \frac{\hat{\theta}-\theta}{\sqrt{\frac{1}{n_1+n_2+n_3-3}(n_1+n_2+n_3 - 3)S_p^2/\sigma^2}\left(\sigma\sqrt{\frac{a_1^2}{n_1}+\frac{a_2^2}{n_2}+\frac{a_3^2}{n_3}}\right)} \\ 
&= \left(\frac{\hat{\theta}-\theta}{\sigma\sqrt{\frac{a_1^2}{n_1}+\frac{a_2^2}{n_2}+\frac{a_3^2}{n_3}}}\right)\frac{1}{\sqrt{\frac{1}{n_1+n_2+n_3-3}(n_1+n_2+n_3 - 3)S_p^2/\sigma^2}}.
\end{align}
Note that $Z=\frac{\hat{\theta}-\theta}{\sigma\sqrt{\frac{a_1^2}{n_1}+\frac{a_2^2}{n_2}+\frac{a_3^2}{n_3}}}$ is a standard normal random variable and $X = (n_1+n_2+n_3 - 3)S_p^2/\sigma^2$ is a chi-squared random variable with $n_1+n_2+n_3-3$ degrees of freedom. Thus, we can write $T$ as
\begin{align}
T &= \frac{Z}{\sqrt{X/(n_1+n_2+n_3-3)}}
\end{align}
If $Z$ and $X$ are independent, then $T$ has a $t$ distribution with $n_1+n_2+n_3-3$ degrees of freedom. To show that $Z$ and $X$ are independent, we only need to show that $\hat{\theta}$ and $S_p^2$ are independent. Now, from Theorem 5.3.1 (Casella and Berger $\S$ 5.3) we know that $\bar{X}$ is independent of $S_1^2$, $\bar{Y}$ is independent of $S_2^2$, and $\bar{W}$ is independent of $S_3^2$. Since $\hat{\theta}$ is a linear combination of $\bar{X}, \bar{Y},$ and $\bar{W}$ and $S_p^2$ is a linear combination of $S_1^2, S_2^2,$ and $S_3^2$, it follows that $\hat{\theta}$ and $S_p^2$ are independent. Therefore, $Z$ and $X$ are independent. Thus, $T\sim t(n_1+n_2+n_3-3)$. 

\vspace{3mm}

\noindent \textcolor{blue}{\textbf{(d)}} Since the variances of the distributions are unknown, the confidence interval for $\theta$ with confidence coefficient $1-\alpha$ is $$\hat{\theta} \pm t_{\alpha/2; n_1+n_2+n_3 - 3}S_p\sqrt{\frac{a_1^2}{n_1} + \frac{a_2^2}{n_2}+\frac{a_3^2}{n_3}}.$$ 

\vspace{3mm}

\noindent \textcolor{blue}{\textbf{(e)}} The hypotheses $H_0: \theta = \theta_0$ versus $H_a: \theta \neq \theta_0$ can be tested using the $T$ statistic from the previous problem. Since the alternative hypothesis contains a ''not equals`` sign, the test will be two-sides with rejection region $$|T| >t_{\alpha/2;n_1+n_2+n_3-3}.$$

\hfill$\blacksquare$

\vfill
\begin{center}
\pgfornament[width = 75mm,
             color = black]{89}
\end{center}
\begin{center}
\vspace{-0.6cm}
\hrulefill \\
\vspace{-0.7cm}
\hrulefill
\end{center}

\pagebreak

\noindent\mybox{\colorbox{lightgray}{\textbf{Wackerly 10.127}}}\hspace{3mm}\hrulefill 

\noindent A merchant figures her weekly profit to be a function of three variables: retail sales (denoted by $X$), wholesale sales (denoted by $Y$), and overhead costs (denoted by $W$). The variables $X$, $Y$, and $W$ are regarded as independent, normally distributed random variables with means $\mu_1, \mu_2,$ and $\mu_3$ and variances $\sigma^2, a\sigma^2,$ and $b\sigma^2$, respectively, for known constants $a$ and $b$ but unknown $\sigma^2$. The merchant's expected profit per week is $\mu_1+\mu_2-\mu_3$. If the merchant has made independent observations of $X$, $Y$, and $W$ for the past $n$ weeks, construct a test of $H_0: \mu_1 + \mu_2 - \mu_3 = k$ against the alternative $H_a: \mu_1 + \mu_2 - \mu_3 \neq k$, for a given constant $k$. You may specify $\alpha = 0.05$. 

\noindent\mybox{\colorbox{lightgray}{\textbf{Solution}}} \hspace{3mm}\hrulefill

\noindent Let $\theta = \mu_1 + \mu_2 - \mu_3$. We can rewrite the hypotheses as $H_0: \theta = k$ versus $H_a: \theta \neq k$. The MLE of $\theta$ is $\hat{\theta} = \bar{X} + \bar{Y} - \bar{W}$, which has a normal distribution with mean $\mu_1 + \mu_2 - \mu_3$ and variance $\frac{\sigma^2}{n}(1+a+b)$. The distribution of $\hat{\theta} - \theta$ is normal with mean $0$ and variance $\frac{\sigma^2}{n}(1+a+b)$. Therefore, the distribution of $$Z = \frac{\hat{\theta}-\theta}{\sigma\sqrt{\frac{1}{n}(1+a+b)}}$$ is standard normal. Next, define $$S_p^2 = \frac{(n-1)S_x^2 + \frac{(n-1)}{a}S_y^2+\frac{n-1}{b}S_w^2}{3n-3},$$ where $S_x^2, S_y^2,$ and $S_w^2$ are the respective sample variances. Then 
\begin{align}
\chi &= (3n-3)S_p^2/\sigma^2 \\
&= \frac{(n-1)S_x^2}{\sigma^2} + \frac{(n-1)S_y^2}{a\sigma^2} + \frac{(n-1)S_w^2}{b\sigma^2}
\end{align}
has a chi-squared distribution with $3n-3$ degrees of freedom. By applying Theorem 5.3.1 (Casella and Berger $\S$ 5.3), we conclude that $Z$ and $\chi$ are independent. So the distribution of 
\begin{align}
\frac{Z}{\sqrt{\chi/(3n-3)}} &= \left(\frac{1}{S_p/\sigma} \right)\left(\frac{\hat{\theta}-\theta}{\sigma\sqrt{\frac{1}{n}(1+a+b)}}\right) \\
&= \frac{\hat{\theta} - \theta}{S_p\sqrt{\frac{1}{n}(1+a+b)}}
\end{align}
is $t$ with $3n-3$ degrees of freedom. So we reject $H_0$ if $$\left|\frac{\hat{\theta} - \theta}{S_p\sqrt{\frac{1}{n}(1+a+b)}}\right| > t_{0.025; 3n-3}.$$

\hfill$\blacksquare$

\vfill
\begin{center}
\pgfornament[width = 75mm,
             color = black]{89}
\end{center}
\begin{center}
\vspace{-0.6cm}
\hrulefill \\
\vspace{-0.7cm}
\hrulefill
\end{center}

\pagebreak

\noindent\mybox{\colorbox{lightgray}{\textbf{Wackerly 10.128}}}\hspace{3mm}\hrulefill 

\noindent A reading exam is given to the sixth graders at three large elementary schools. The scores on the exam at each school are regarded as having normal distributions with unknown means $\mu_1, \mu_2,$ and $\mu_3$, respectively, and unknown common variance $\sigma^2$ ($\sigma_1^2 = \sigma_2^2 = \sigma_3^2 = \sigma^2$). Using the data in the accompanying table on independent random samples from each school, test to see if evidence exists of a difference between $\mu_1$ and $\mu_2$. Use $\alpha = 0.05$.

\noindent\mybox{\colorbox{lightgray}{\textbf{Solution}}} \hspace{3mm}\hrulefill

\noindent This problem can be solved using the theory developed in problem 10.126. Letting $\theta = \mu_1 - \mu_2 + 0\mu_3 = \mu_1-\mu_2$, our hypothesis test becomes $H_0: \theta = 0$ versus $H_1: \theta \neq 0$. We will reject $H_0$ if the test statistic $$T = \frac{\bar{X}-\bar{Y}-(\mu_1-\mu_2)}{S_p\sqrt{\frac{1}{10}+\frac{1}{10}}}$$ satisfies $|T| > t_{0.025; 27} = 2.051831$. Using the data from the table, we find that
\begin{align}
\bar{x} - \bar{y} &= 60 - 50 \\
&= 10,
\end{align}
and
\begin{align}
S_p^2 &= \frac{9(S_x^2+S_y^2+S_w^2)}{27} \\
&= \frac{9}{27}\left( \frac{\sum x_i^2 - n\bar{x}^2}{9} + \frac{\sum y_i^2 - n\bar{y}^2}{9} + \frac{\sum w_i^2 - n\bar{w}^2}{9}\right) \\
&= \frac{1}{27}[36950 - 10(60^2) + 25850 - 10(50^2) + 49900 - 10(70^2)] \\
&= \frac{2700}{27} \\
&= 100.
\end{align}
Thus, 
\begin{align}
T &= \frac{10}{10\sqrt{2/10}} \\
&= \sqrt{5} \\
& \approx 2.23607.
\end{align}
Since $|T| = 2.23607 > 2.051831$, we reject $H_0$ and conclude that there is evidence supporting the hypothesis that there is a difference between $\mu_1$ and $\mu_2$. 

\hfill$\blacksquare$

\vfill
\begin{center}
\pgfornament[width = 75mm,
             color = black]{89}
\end{center}
\begin{center}
\vspace{-0.6cm}
\hrulefill \\
\vspace{-0.7cm}
\hrulefill
\end{center}

\pagebreak

\noindent\mybox{\colorbox{lightgray}{\textbf{Wackerly 10.129}}}\hspace{3mm}\hrulefill 

\noindent Suppose that $Y_1, Y_2, \ldots, Y_n$ denote a random sample from the probability density function given by 
$$
f(y|\theta_1, \theta_2) = 
\begin{cases}
\frac{1}{\theta_1}e^{-(y-\theta_2)/\theta_1}, & y > \theta_2 \\
0, & \text{elsewhere.}
\end{cases}
$$
Find the likelihood ratio test for testing $H_0:\theta_1 = \theta_{1, 0}$ versus $H_a: \theta_1 > \theta_{1, 0}$ with $\theta_2$ unknown.

\noindent\mybox{\colorbox{lightgray}{\textbf{Solution}}} \hspace{3mm}\hrulefill

\noindent The likelihood function for the sample is
\begin{align}
L(\theta_1, \theta_2|\textbf{y}) &= \prod_{i=1}^n f(y_i|\theta_1, \theta_2) \\
&= \begin{cases}
\left(\frac{1}{\theta_1}\right)^n \exp\left\lbrace -\frac{1}{\theta_1}\sum\limits_{i=1}^n(y_i-\theta_2)\right\rbrace, & \theta_2 \leq y_{(1)} \\
0, & \text{elsewhere}
\end{cases}
\end{align}
The log-likelihood function for $\theta_2 < y_{(1)}$ is then 
\begin{align}
\mathscr{L}(\theta_1,\theta_2|\textbf{y}) &= -n\log(\theta_1)-\frac{1}{\theta_1}\sum\limits_{i=1}^n(y_i - \theta_2).
\end{align}
For other values of $\theta_2$, it doesn't exist. Now we will use these two functions to find the MLEs of $\theta_1$ and $\theta_2$ (which we will denote by $\hat{\theta}_1$ and $\hat{\theta}_2$, respectively). Examining the likelihood function, we see that on the interval $-\infty < \theta_2 < y_{(1)}$, it is an increasing function of $\theta_2$. Thus, the MLE of $\theta_2$ is $\hat{\theta}_2 = y_{(1)}$. To find the MLE of $\theta_1$, we solve
\begin{align}
\frac{d}{d\theta_1}\mathscr{L}(\theta_1, \hat{\theta}_2|\textbf{y}) &= -\frac{n}{\theta_1}+\frac{1}{\theta_1^2}\sum\limits_{i=1}^n(y_i-y_{(1)}) \\
&= 0
\end{align}
for $\theta_1$.  It is easy to see that the solution is $\hat{\theta}_1 = \frac{1}{n}\sum\limits_{i=1}^n (y_i - y_{(1)}).$

The likelihood ratio test statistic for the hypotheses in this problem is 
\begin{align}
\lambda_1(\textbf{y}) &= \frac{L(\theta_{1,0}, \hat{\theta}_2|\textbf{y})}{L(\hat{\theta}_1, \hat{\theta}_2|\textbf{y})} \\
&= \left(\frac{\hat{\theta}_1}{\theta_{1,0}}\right)^n \frac{\exp\left\lbrace -\frac{1}{\theta_{1,0}}\sum(y_i-y_{(1)})\right\rbrace}{\exp\left\lbrace -\frac{1}{\hat{\theta}_1}\sum(y_i-y_{(1)})\right\rbrace} \\
&= \left(\frac{\sum(y_i-y_{(1)})}{n\theta_{1,0}}\right)^n\exp\left\lbrace -\frac{1}{\theta_{1,0}}\sum(y_i-y_{(1)}) + n\right\rbrace.
\end{align}
The likelihood ratio test is to reject $H_0$ whenever $\lambda_1(\textbf{y}) < c$, where $0\leq c \leq 1$. 

\hfill$\blacksquare$

\vfill
\begin{center}
\pgfornament[width = 75mm,
             color = black]{89}
\end{center}
\begin{center}
\vspace{-0.6cm}
\hrulefill \\
\vspace{-0.7cm}
\hrulefill
\end{center}

\pagebreak

\noindent\mybox{\colorbox{lightgray}{\textbf{Wackerly 10.130}}}\hspace{3mm}\hrulefill 

\noindent Refer to Exercise 10.129. Find the likelihood ratio test for testing $H_0: \theta_2 = \theta_{2,0}$ versus $H_a: \theta_2 > \theta_{2,0}$, with $\theta_1$ unknown. 

\noindent\mybox{\colorbox{lightgray}{\textbf{Solution}}} \hspace{3mm}\hrulefill

\noindent If $\theta_1$ and $\theta_2$ are restricted to values specified by the null hypothesis, then $L$ is maximized when $\theta_2 = \theta_{2,0}$ and $\theta_1 = \frac{1}{n}\sum(y_i-\theta_{2,0})$. Thus, the likelihood test statistic for this hypothesis test is
\begin{align}
\lambda_2(\textbf{y}) &= \frac{L\left(\frac{1}{n}\sum(y_i - \theta_{2,0}), \theta_{2,0} | \textbf{y}\right)}{L(\hat{\theta}_1, \hat{\theta}_2|\textbf{y})} \\
&= \left[\frac{\frac{1}{n}\sum(y_i-y_{(1)})}{\frac{1}{n}\sum(y_i-\theta_{2,0})}\right]^n \frac{\exp\left\lbrace -\frac{\sum(y_i-\theta_{2,0})}{\frac{1}{n}\sum(y_i-\theta_{2,0})}\right\rbrace}{\exp\left\lbrace -\frac{\sum(y_i-y_{(1)})}{\frac{1}{n}\sum(y_i-y_{(1)})}\right\rbrace} \\
&= \left[\frac{\sum(y_i-y_{(1)})}{\sum(y_i-\theta_{2,0})}\right]^n \frac{\exp\left\lbrace-n\right\rbrace}{\exp\left\lbrace-n\right\rbrace} \\
&= \left[\frac{\sum(y_i-y_{(1)})}{\sum(y_i-\theta_{2,0})}\right]^n. 
\end{align}
Again, the likelihood ratio test is to reject $H_0$ whenever $\lambda_1(\textbf{y}) < c$, where $0\leq c \leq 1$.

\hfill$\blacksquare$

\vfill
\begin{center}
\pgfornament[width = 75mm,
             color = black]{89}
\end{center}
\begin{center}
\vspace{-0.6cm}
\hrulefill \\
\vspace{-0.7cm}
\hrulefill
\end{center}

\pagebreak

\end{document}
